\section{Introduction}

Les données de santé sont relationnelles, mais la présence de liens ne garantit pas que leur modélisation améliore la prédiction. Pour mesurer leur valeur, il faut distinguer l'information contenue dans les relations, les méthodes qui l'exploitent et la cible évaluée.

HealthGraphBench met cette question à l'épreuve sur trois tâches de santé réglementaire : prédire de nouveaux liens produit--problème dans MAUDE, des déficiences d'inspection dans les maisons de retraite suivies par CMS et de nouveaux liens prescripteur--médicament dans Medicare Part D. Le benchmark réunit définitions des tâches, périodes, ensembles de candidats, références interprétables et résultats historiques ; les comparaisons Q2 ajoutent des familles de modèles nouvellement réglées sur chacun des trois domaines, en plus des analyses historiques C et D1 sur MAUDE.

L'étude compare les performances de classement face à des références simples ; elle n'évalue ni l'utilité clinique ni les décisions qui pourraient découler des scores.

\section{Travaux connexes : des performances aux conditions de comparaison}

\subsection{Périmètre de la revue}

La revue narrative, arrêtée au 29 septembre 2026, porte sur les benchmarks relationnels, l'évaluation temporelle, les références simples et la validité des études d'apprentissage automatique. Elle s'appuie sur des articles originaux, leurs pages de publication et les documentations officielles. Cette sélection ciblée situe HealthGraphBench dans la littérature sans prétendre en recenser exhaustivement les travaux ni mesurer la fréquence des succès des méthodes relationnelles.

\subsection{Benchmarks généraux et données relationnelles}

Open Graph Benchmark a structuré l'évaluation autour de jeux de données, partitions et métriques partagés\cite{hu2020ogb}. RelBench déplace l'attention vers la prédiction sur des bases relationnelles, où plusieurs tables doivent être exploitées conjointement\cite{robinson2024relbench}. RelBench v2 étend cette approche et introduit notamment des tâches de complétion respectant des contraintes temporelles\cite{gu2026relbench}. Ces travaux rapportent des bénéfices de l'apprentissage relationnel dans leurs propres conditions expérimentales. Ils constituent un contrepoint important à toute interprétation générale d'un résultat négatif.

HealthGraphBench complète ces ressources par des tâches réglementaires définies, des comparateurs relationnels interprétables et une description des limites de l'inférence. Il ne vise pas à reproduire les classements d'OGB ou de RelBench.

\subsection{Temporalité, candidats et références simples}

Temporal Graph Benchmark montre que les performances varient fortement selon les jeux de données et que certaines méthodes simples restent compétitives\cite{huang2023tgb}. Poursafaei et ses collègues soulignent également le rôle des stratégies d'évaluation et proposent une référence fondée sur la mémorisation des liens\cite{poursafaei2022}. Ces enseignements invitent à examiner l'univers des candidats avant d'interpréter un score.

La transposition à HealthGraphBench exige toutefois une précaution : MAUDE et Part D excluent les relations déjà observées pour l'entité cible. La seule répétition d'un lien mémorisé ne peut donc expliquer leurs résultats. En revanche, la popularité d'une catégorie ou les observations de pairs peuvent encore être prédictives. Évaluer tous les candidats admissibles évite que la difficulté soit déterminée par un petit échantillon de négatifs faciles ; cela ne rend pas les absences observées équivalentes à de véritables négatifs cliniques.

Les exigences de comparaison équitable étudiées par Errica et ses collègues rappellent enfin que l'architecture ne peut être dissociée du protocole et des références\cite{errica2020}. Leurs expériences portent sur la classification de graphes, non sur les tâches présentes : nous en retenons un principe méthodologique, pas une comparaison numérique directe.

\subsection{Validité scientifique et transparence}

Les travaux sur les fuites d'information et le cadre REFORMS motivent la description explicite des cibles, unités, données, partitions et comparateurs\cite{kapoor2023leakage,kapoor2024reforms}. Ils éclairent le rapport de l'étude ; ils ne constituent pas une certification de conformité à REFORMS.

Un ordre chronologique correct est nécessaire, mais insuffisant. Il ne garantit ni la disponibilité publique historique de chaque champ, ni l'absence d'adaptation des choix après consultation des résultats. Ces deux limites sont particulièrement importantes pour des registres dont les observations peuvent être publiées ou corrigées après l'événement décrit.


\section{Méthodes}
\label{sec:methods}

\subsection{Question et cadre des comparaisons}

Les résultats historiques de MAUDE et CMS proviennent du socle v0.1 ; Part D a été ajoutée dans la version v0.2. Les analyses de durée C et de suppression de voisinage D1 sur MAUDE sont également conservées comme résultats historiques distincts. Elles ne sont ni remplacées ni réentraînées par les comparaisons Q2, qui ajoutent de nouveaux ajustements sur les trois domaines.

Les périodes de validation et de test utilisées pour ces travaux avaient déjà été consultées au cours du développement antérieur. Les sélections Q2 sont néanmoins écrites à partir de la validation avant l'évaluation des tests, sans choix fondé sur les métriques de test. Cela ne rend pas les périodes indépendantes : les résultats Q2 demeurent exploratoires et n'apportent pas de confirmation indépendante.

Les expériences historiques et leurs sorties sont décrites dans les rapports et le dépôt du benchmark\cite{hgb2026,hgbcontract,hgbresults,hgbpartd,hgbanalysis,hgbcode}. Le DOI \nolinkurl{10.5281/zenodo.22796551} désigne la version 0.2 du benchmark. Les ajustements Q2 et leurs preuves distinctes sont décrits dans les sections S12--S14\cite{hgbq2methods}.

L'analyse C compare des durées GraphSAGE puis rapporte le test de la durée choisie ; D1 confronte le modèle à agrégation moyenne à un contrôle sans propagation sous réglages imposés. Ces analyses restent séparées des familles Q2 réglées ci-dessous.

\subsection{Tâches, populations et disponibilité des informations}
\label{sec:tasks}

\paragraph{MAUDE.} Pour un code produit, prédire les relations avec des codes problèmes observées pour la première fois dans la fenêtre conservée. L'horloge est \texttt{DATE\_RECEIVED}, non une date vérifiée de publication de chaque champ. Un produit doit avoir au moins un signalement antérieur ; un problème doit déjà exister globalement ; la relation doit être absente de l'historique du produit. Le « seuil 1 » est donc un seuil de soutien historique du produit en nombre de signalements, non un seuil de gravité. Le classement porte sur tous les problèmes admissibles. Les nouveautés hors vocabulaire antérieur ne font pas partie de ce problème de classement\cite{hgbcontract,hgbcode}.

\paragraph{CMS nursing homes.} Pour les inspections Health Standard retenues de maisons de retraite médicalisées américaines, prédire la présence d'une déficience standard G--L constatée à l'inspection cible à partir d'enregistrements datés d'avant celle-ci. Ce filtre de dates ne prouve pas leur disponibilité publique historique. L'unité est un épisode (CCN, date), avec au moins une inspection antérieure. Le score est construit rétrospectivement juste avant cet épisode : la tâche est conditionnée à la tenue de l'inspection. Elle ne prédit ni sa survenue ni sa date, et ne correspond pas à un classement de tous les établissements au début de l'année. Les extraits de fournisseurs actuels et le faible nombre de cycles retenus induisent sélection et troncature historique\cite{cmsnursing,hgbcontract,hgbcode}.

\paragraph{Medicare Part D.} Le médicament est le nom générique exact, débarrassé des espaces périphériques ; il n'est pas normalisé en concept RxNorm. Chaque cohorte cible comprend au plus 2\,000 prescripteurs déjà observés, sélectionnés par l'ordre (SHA-256 du NPI, NPI). Tous les médicaments présents dans l'historique antérieur de la cohorte, sauf ceux déjà associés au prescripteur, sont classés. Le positif est un nouveau lien \emph{publié et observé}, pas une première prescription. Le contrat rappelle la suppression CMS des combinaisons ayant au plus dix demandes : une absence n'est pas un négatif clinique\cite{cmspartd,hgbpartd}.

\begin{table}[!htbp]
\centering\footnotesize
\caption{Effectifs récupérables pour les partitions évaluées. Les observations MAUDE sont les produit--trimestres avec positif ; celles de CMS, les inspections cibles retenues ; celles de Part D, les prescripteur--années admissibles. Les unités utilisées pour calculer les intervalles ne représentent pas les effectifs nationaux. Sources : rapports, audit des dénominateurs, analyse annuelle et extraction des descripteurs\cite{hgbresults,hgbanalysis,hgbpartd}; détails en S3.}
\label{tab:population}
\begin{tabularx}{\linewidth}{@{}lrrX@{}}
\toprule
Tâche / partition & Observations & Positifs & Unités distinctes documentées\\
\midrule
MAUDE, validation 2023 & 3\,483 & 7\,705 liens & 1\,641 produits\\
MAUDE, test 2024--2025 & 6\,370 & 13\,174 liens & 2\,026 produits, grappes\\
CMS, validation 2023 & 2\,661 & 408 & 2\,660 CCN\\
CMS, test 2024--2025 & 18\,985 & 2\,423 & 13\,889 CCN, grappes\\
Part D, validation 2023 & 2\,000 & 5\,794 liens & 1\,035 avec positif\\
Part D, test 2024 & 2\,000 & 5\,476 liens & 1\,019 avec positif\\
\bottomrule
\end{tabularx}
\end{table}

Le supplément S3 donne les effectifs historiques, d'entraînement et de candidats Part D. Pour CMS, les constantes de prévalence et le protocole à fenêtre croissante permettent de reconstituer 2\,066, 4\,727 et 13\,636 lignes d'entraînement avant 2023, 2024 et 2025, dont 244, 652 et 1\,905 positives; cette reconstruction n'est pas une lecture directe des lignes.

Les entrées tabulaires MAUDE comptent 96\,906 à 124\,322 lignes aux trois réajustements (48\,453 à 62\,161 positives, avec un non-positif échantillonné par positif); les graphes ont 67\,676 à 81\,891 arêtes historiques. Les boucles gelées impliquent 270\,704 à 327\,564 pas de gradient BPR et 203\,028 à 245\,673 pas GraphSAGE, non des journaux d'optimiseur récupérés.

Dans les inventaires reconstruits pour l'analyse historique B1/B2, 97,56\,\% des paires candidates et 97,80\,\% des liens positifs sont couvertes au test 2024--2025. Pourtant, 1\,414 des 6\,370 produit--trimestres (22,20\,\%) ont tous leurs candidats admissibles représentés. Le supplément S10 détaille ces comptes. Cet audit historique décrit la couverture des identifiants reconstruits et n'isole pas l'effet du repli à zéro ; contrairement à ces reconstructions, les nouveaux checkpoints d'entraînement sont inspectés en S11.

\begin{table}[!htbp]
\centering\footnotesize
\caption{Périodes et unités d'évaluation des trois tâches. L'ordre des événements ne garantit pas la disponibilité publique historique de tous les attributs.}
\label{tab:time}
\begin{tabularx}{\linewidth}{@{}lXXX@{}}
\toprule
 & MAUDE & CMS nursing homes & Part D\\
\midrule
Historique initial & 2019T1--2022T4 & 2019--2022 & 2019--2022\\
Validation / test & 2023 / 2024--2025 & 2023 / 2024--2025 & 2023 / 2024\\
Instant du score & Avant le trimestre cible, selon l'horloge des signalements & Juste avant l'inspection cible retenue & Avant l'année de service cible, en reconstruction rétrospective\\
Mise à jour & Après chaque trimestre évalué ; réajustement annuel & Entraînement sur les années strictement antérieures & Historique 2019--2023 pour le test 2024\\
Horizon / cible & Liens du trimestre & Déficience de cet épisode, sans délai de déploiement validé & Liens publiés de l'année de service\\
\bottomrule
\end{tabularx}
\end{table}

L'évaluation MAUDE est roulante : après chaque trimestre, les nouvelles relations enrichissent l'historique utilisé pour le trimestre suivant ; les réajustements annuels ont lieu au début de 2023, 2024 et 2025. Cette mise à jour prévue par le protocole est distincte de l'adaptation des choix après consultation des résultats. Pour CMS, une date d'association de propriété antérieure à l'inspection ne garantit pas que cette association était alors publiquement connue ; pour Part D, l'année de service ne correspond pas à la date de publication du fichier\cite{hgbcode,hgbresults}.

\subsection{Heuristiques et modèles effectivement comparés}
\label{sec:models}

\paragraph{MAUDE : fréquence des voisins.} Soit $H_p$ l'ensemble des problèmes antérieurement associés au produit $p$, et $U_d$ l'ensemble des produits ayant déjà présenté le problème $d$. Les voisins sont $N_p=(\bigcup_{d\in H_p} U_d)\setminus\{p\}$. Pour un candidat $c\notin H_p$, les deux scores sont
\begin{equation}
 s_{\mathrm{global}}(p,c)=|U_c|,\qquad
 s_{\mathrm{voisins}}(p,c)=|N_p\cap U_c|.
\end{equation}
Chaque voisin distinct compte une fois : ni pondération par volume de signalements ni normalisation par degré n'est appliquée. Sans soutien, le score vaut zéro. Les égalités sont départagées par soutien global décroissant, puis code problème croissant\cite{hgbcode}. Cette définition précise ce que personnalise l'heuristique.

Les autres comparateurs historiques MAUDE sont une régression logistique, des arbres à un niveau boostés, une factorisation spectrale de rang 8, un BPR avec moyenne fixe à un saut et le GraphSAGE historique à un saut\cite{hamilton2017,rendle2009,hgbcode}. Celui-ci combine des vecteurs de nœuds entraînables et une moyenne de vecteurs voisins au moyen de deux transformations partagées, d'une activation $\tanh$ et d'un score par produit scalaire. Le modèle historique utilise la dimension 8, huit voisins, trois époques, un taux d'apprentissage de 0,02 et une régularisation de 0,0005. L'analyse de durée et le contrôle sans agrégation sont décrits dans les résultats et en S11. La généralisation aux nœuds absents du réajustement n'a pas été évaluée séparément ; ces cas reçoivent un score de repli nul, et leur fréquence est rapportée en S10.

\paragraph{Part D : spécialité et recouvrement.} La popularité par spécialité compte les prescripteurs historiques de la même spécialité associés au médicament. La spécialité est celle de la dernière année observée, si elle est unique et non vide ; sinon, le score revient au soutien global. Pour les historiques $D_i,D_j$ des prescripteurs, le recouvrement utilise
\begin{equation}
 s_{\mathrm{recouvrement}}(i,d)=\sum_{j\ne i}\frac{|D_i\cap D_j|}{|D_i\cup D_j|}\,\mathbf{1}\{d\in D_j\}.
\end{equation}
Seuls les pairs ayant une intersection non vide contribuent ; la somme n'est pas normalisée par le nombre de pairs. Les scores nuls suivent la règle commune de départage : soutien global décroissant puis clé canonique croissante. Les modèles appris sont la logistique sur dix variables historiques et un BPR prescripteur--médicament enrichi par une représentation de spécialité. Ce BPR n'effectue pas de propagation de messages\cite{hgbcode,hgbpartd}.

\paragraph{CMS : augmentation par propriétaires.} Les deux modèles historiques sont des régressions logistiques : onze variables numériques locales et des indicatrices d'État, puis douze agrégats relationnels supplémentaires. La représentation \emph{combinée} privilégie les identifiants de propriétaires CHOW ; les noms normalisés courants servent de repli pour les CCN absents de cette représentation. Il ne s'agit pas d'une union systématique des deux ensembles. Les pairs sont les autres CCN partageant une association dont la date de début est au plus celle de la cible. Les dates de fin n'étant pas toutes connues, cette règle n'atteste pas une propriété juridiquement active. Les nombres d'inspections et de déficiences des pairs sont cumulés strictement avant la cible ; les taux divisent les totaux de déficiences par les totaux d'inspections, avec zéro en l'absence de dénominateur\cite{hgbcode}.

Les deux modèles CMS historiques utilisent 300 itérations, un taux d'apprentissage de 0,08 et une pénalité $L_2=0,05$, après standardisation sur les données d'entraînement. Le supplément décrit les paramètres des autres méthodes. Ces paramètres sont ceux des exécutions conservées, non le résultat d'une recherche exhaustive ou d'une comparaison des coûts.
 
\subsection{Sélection des nouvelles comparaisons Q2}
\label{sec:q2-methods}
La campagne complémentaire porte l'identifiant Q2 dans les artefacts ; ce nom distingue ses nouveaux ajustements des configurations historiques.

Chaque famille réglable comprend six configurations publiées ; les heuristiques ne sont pas réglées. La configuration est sélectionnée sur la moyenne du rappel micro@10 de validation pour MAUDE et Part D, et sur l'AUC ROC de validation pour CMS, avec l'ordre de grille comme départage. Les modèles stochastiques utilisent les graines 103, 211 et 307 ; la sélection porte sur leur moyenne, jamais sur la meilleure graine. Les HGB sont ajustés trois fois seulement lorsque l'entraînement exact dépasse 200\,000 lignes ; sinon, une seule exécution est faite.

Pour les GraphSAGE \texttt{mean} et \texttt{none} sur MAUDE, chacune des trois graines doit modifier l'état du modèle et réduire d'au moins 1\,\% la perte sur la même sonde fixe construite à partir du seul historique. Seules les configurations admissibles participent à la sélection et les réajustements retenus repassent ces critères avant l'évaluation. \texttt{none} réalise $\tanh(W_{\mathrm{self}}x)$, sans propagation de voisinage, mais ses vecteurs de nœuds sont toujours appris sur les liens historiques avec BPR. Le modèle \texttt{mean} active deux transformations de dimension $d$ (soit $2d^2$ paramètres scalaires) contre $d^2$ pour \texttt{none}, en plus des vecteurs de nœuds : les capacités ne sont pas égalisées et le contraste n'identifie pas un effet causal de l'agrégation. Le détail du protocole est en S12.

Les six opportunités de sélection ne correspondent pas à un coût de calcul identique : dimensions, répétitions et durée des ajustements diffèrent, et les exécutions ont partagé la machine. Les ressources observées sont rapportées avec les limites d'interprétation en S12.

\subsection{Mesures, égalités et incertitude}
\label{sec:metrics}

Pour une observation $i$ comportant $P_i>0$ positifs et $h_i(K)$ liens retrouvés, le rappel micro est $\sum_i h_i(K)/\sum_i P_i$ et le rappel macro est la moyenne des $h_i(K)/P_i$. La MRR moyenne l'inverse du rang du premier positif. Pour Part D, précision et charge incluent aussi les observations sans positif : avec $a_i(K)=\min(K,|C_i|)$, la précision est $\sum_i h_i(K)/\sum_i a_i(K)$ et la charge moyenne est $\sum_i a_i(K)/n$. Un seuil de classement n'est pas une confiance calibrée.

CMS rapporte l'AUC ROC, le Brier, la précision et le rappel au décile, ainsi que deux conventions distinctes de précision moyenne. L'\emph{AP après départage historique}, notée $\mathrm{AP}_{\mathrm{rang}}$, moyenne les précisions aux rangs positifs après tri par score décroissant, puis indice de ligne croissant (date puis CCN). L'\emph{AP par seuils}, notée $\mathrm{AP}_{\mathrm{seuils}}$, regroupe les scores exactement égaux. Pour un groupe de seuil $j$ contenant $t_j$ positifs, avec $T_j$ positifs parmi les $N_j$ observations cumulées,
\begin{equation}
 \mathrm{AP}_{\mathrm{seuils}}=\sum_j\frac{t_j}{P}\frac{T_j}{N_j}.
\end{equation}
Cette convention est une moyenne non interpolée des précisions aux seuils\cite{sklearnap}. L'AUC traite les égalités par rang moyen ; le décile retient les $\lceil0,1n\rceil$ premières lignes. Les résultats et intervalles historiques d'AP concernent uniquement $\mathrm{AP}_{\mathrm{rang}}$ ; aucun intervalle n'est associé à la nouvelle AP par seuils (détails en S7).

Les intervalles historiques au v0.1 sont signalés comme tels au tableau~\ref{tab:ci}. Deux bootstraps appariés post hoc complètent ces résultats sur les prédictions conservées : par produit pour MAUDE et par prescripteur pour Part D (S6). Un nouveau rapprochement GraphSAGE30--voisins apparie les contributions produit--trimestre et rééchantillonne 2\,026 produits, en conservant leurs trimestres : 1\,000 tirages, graine 20261003, IC percentile à 95\,\% avec interpolation linéaire et rapport regroupé recalculé à chaque tirage (S11). Ces intervalles sont conditionnels aux prédictions et à la population observées ; ils ne couvrent ni la sélection, ni l'entraînement, ni toute dépendance entre produits. Aucun nouvel IC n'est calculé pour la durée elle-même ou D1. Les comparaisons Q2 n'ajoutent aucun bootstrap ni intervalle : leurs résultats sont des estimations ponctuelles.


\section{Résultats}

\subsection{MAUDE historique : avantage à dix propositions, dépendance au seuil}

\begin{table}[!htbp]
\centering\small
\caption{Résultats historiques MAUDE v0.1, test 2024T1--2025T4, soutien historique du produit $\geq1$. GraphSAGE désigne ici la configuration historique à trois époques, non GraphSAGE30. R@K : rappel micro. Le gras indique le meilleur point parmi les modèles historiques de chaque colonne, sans test supplémentaire de supériorité. Source : tableau unifié\cite{hgbresults}.}
\label{tab:maude}
\input{tables/maude_fr.tex}
\end{table}

Dans les résultats historiques, la fréquence des voisins atteint un rappel à 10 de 0,192045, contre 0,184454 pour la popularité globale et 0,172840 pour GraphSAGE à trois époques (tableau~\ref{tab:maude}). L'IC post hoc par produit pour voisins moins global est $[+0,005497;+0,009820]$, pour une différence de $+0,007591$. L'intervalle publié pour GraphSAGE moins voisins concerne lui aussi les seules prédictions historiques : aucun de ces intervalles ne décrit les nouvelles analyses. Le rappel macro et la MRR suivent le même classement historique\cite{hgbresults}.

L'heuristique ne domine cependant pas toutes les métriques. À $K=20$, la factorisation spectrale atteint 0,293912, contre 0,293077 pour la fréquence des voisins : l'écart descriptif est de 0,000835 (figure~\ref{fig:cutoffs}). Aucun intervalle n'est disponible ici pour ce contraste. Ce changement de classement, même faible, interdit une conclusion uniforme sur tous les budgets de sélection. Il ne remet pas en cause le contraste principal à $K=10$.

\begin{figure}[!ht]
\centering
\includegraphics[width=\linewidth]{figures/maude_cutoffs_historical3_fr.pdf}
\caption{Différences de rappel MAUDE par rapport aux voisins. À $K=10$, l'IC de GraphSAGE historique à trois époques moins voisins et l'IC complémentaire de global moins voisins sont distingués dans la légende ; aucun ne concerne GraphSAGE30. L'IC global moins voisins vaut $[-0,009820;-0,005497]$ : son signe et l'ordre de ses bornes sont inversés par rapport à voisins--global (tableau~\ref{tab:ci}). La quasi-égalité à $K=20$ concerne factorisation spectrale et voisins.}
\label{fig:cutoffs}
\end{figure}

\subsection{MAUDE : exploration complémentaire de la durée}
\label{sec:maude-post-rc1}

La grille de validation 0/3/10/30 retient 30 époques sur 2023 (0,205062). Sa sensibilité n'est pas monotone : dix époques font moins bien que trois (figure~\ref{fig:duration-main}). Au test regroupé 2024--2025, le modèle retenu atteint 0,210642 (2\,775/13\,174), contre 0,192045 (2\,530/13\,174) pour les voisins conservés. Le rapprochement vérifie les clés produit--trimestre, soutiens et identités positives ; les candidats C concordent avec le même historique préparé et les cardinalités historiques. Les listes historiques complètes n'étant pas exportées, cette dernière vérification combine reconstruction et traces conservées, non une comparaison directe de deux exports complets.

Le contraste apparié GraphSAGE30 moins voisins vaut $+0,018597$, soit environ $+1,86$ point de rappel (tableau~\ref{tab:paired-maude}). L'IC conditionnel par produit est $[+0,011975;+0,025570]$. Il décrit l'incertitude du contraste observé sous les prédictions et la population fixées, sans corriger l'exposition antérieure aux tests ou la sélection de la durée. Ce rapprochement ne mesure pas l'effet causal de passer de trois à trente époques.

\begin{figure}[!htbp]\centering
\includegraphics[width=\linewidth]{figures/maude_duration_RC2_fr.png}
\caption{Exploration de durée MAUDE (C), distincte du benchmark historique : rappel micro à 10 en validation selon des entraînements séparés à 0, 3, 10 et 30 époques, réinitialisés de manière déterministe et sans reprise de l'état d'un entraînement précédent. Il ne s'agit pas de répétitions aléatoires indépendantes. La relation observée n'est pas monotone ; seule la durée retenue sur validation est évaluée dans les tests C. Les pertes d'entraînement sont rapportées en S11.}
\label{fig:duration-main}
\end{figure}

\begin{table}[!htbp]\centering\small
\caption{Rapprochement apparié de GraphSAGE30 sélectionné sur validation et des voisins historiques, test 2024--2025. L'IC est conditionnel à ces prédictions ; les intervalles du GraphSAGE historique à trois époques ne sont pas transférés.}
\label{tab:paired-maude}
\input{tables/maude_paired_RC2_fr.tex}
\end{table}

 
Les critères d'apprentissage des ajustements Q2 sont distincts de D1 et ne réinterprètent pas son résultat historique : la perte presque plate de D1 demeure un fait conservé, mais n'établit pas un bug d'optimisation. Les critères d'admissibilité Q2 ne démontrent ni que D1 apprenait efficacement sous son réglage ni le mécanisme du changement de classement.

Les cohortes évaluées comprennent uniquement les produit--trimestres avec au moins un lien positif ; la précision et la charge sur l'ensemble des observations ne sont pas estimées. Les résultats annuels, dénominateurs et courbes de perte figurent en S11.

\subsection{Inspections : incrément faible et incertain des propriétaires}

\begin{table}[!htbp]
\centering\small
\caption{Inspections CMS, prédictions de test 2024--2025 regroupées. $P_{10\%}$ et $R_{10\%}$ : précision et rappel au décile calculé sur les scores regroupés, non séparément par année. Le Brier doit diminuer ; les autres mesures doivent augmenter. Source : tableau unifié\cite{hgbresults}.}
\label{tab:cms}
\input{tables/cms_fr.tex}
\end{table}

L'historique local atteint une AUC de 0,619975 et l'augmentation par les propriétaires 0,621293 (tableau~\ref{tab:cms}). L'écart de 0,001318 a un intervalle historique qui contient zéro. $\mathrm{AP}_{\mathrm{rang}}$ augmente également très peu, tandis que le Brier et la précision au décile supérieur regroupé se dégradent légèrement en valeur ponctuelle. L'augmentation n'apporte donc pas un bénéfice cohérent sur l'ensemble des mesures\cite{hgbresults}.

L'absence d'un seuil d'équivalence défini interdit de conclure à l'équivalence des deux modèles. Les résultats annuels montrent également que l'écart d'AUC varie de $+0,000545$ en 2024 à $-0,000087$ en 2025 ; la valeur regroupée n'est pas leur moyenne arithmétique\cite{hgbresults}.


\subsection{CMS : discrimination annuelle et regroupement}
\label{sec:cmsannual}

\begin{table}[!htbp]
\centering\footnotesize
\caption{CMS : résultats de test par année, transcrits de l'analyse gelée. Les prévalences sont 14,0644\,\% en 2024 et 11,6118\,\% en 2025. L'AP utilise la convention de rang du code. Le supplément inclut la validation 2023\cite{hgbanalysis}.}
\label{tab:cmsannual}
\input{tables/cms_annual_test_fr.tex}
\end{table}

Les 8\,909 inspections de 2024 comprennent 1\,253 positives ; les 10\,076 inspections de 2025 en comprennent 1\,170. Calculé à partir des AUC non arrondies, l'écart propriétaires moins historique vaut $+0,000545$ en 2024 et $-0,000087$ en 2025, après arrondi à six décimales.

Au décile supérieur, les paires précision/rappel $(P_{10\%},R_{10\%})$ sont en 2024 : prévalence (0,158249 ; 0,112530), historique local (0,273850 ; 0,194733), propriétaires (0,267116 ; 0,189944) ; en 2025 : prévalence (0,123016 ; 0,105983), historique local (0,237103 ; 0,204274), propriétaires (0,235119 ; 0,202564). Le Brier augmente légèrement les deux années.

Les résultats annuels décrivent mieux la discrimination au sein d'une année que le seul résultat regroupé.

Une décomposition descriptive supplémentaire sépare les paires positif--négatif appartenant à la même année des paires inter-annuelles. Si $P_y,N_y$ sont les effectifs positifs et négatifs, $T=(\sum_yP_y)(\sum_yN_y)$ et $W=\sum_yP_yN_y$, alors
\begin{equation}
 A_{\text{regroupée}}=\frac{W}{T} A_{\mathrm{intra}}+
 \frac{T-W}{T} A_{\mathrm{inter}},\qquad
 A_{\mathrm{intra}}=\frac{\sum_y P_yN_y A_y}{W}.
\label{eq:decomp}
\end{equation}
Le terme inter-annuel est reconstitué algébriquement à partir des AUC publiées, sans lecture de nouveaux scores individuels. Les poids sont 0,498707 et 0,501293. L'AUC intra-annuelle est 0,625299 pour l'historique et 0,625515 avec propriétaires : l'écart est $+0,000216$. Les contributions pondérées intra- et inter-annuelles à l'écart regroupé de $+0,001318$ sont respectivement $+0,000108$ et $+0,001210$. Il s'agit d'une identité descriptive conditionnelle aux résumés disponibles, sans nouvel intervalle ni test de supériorité. L'intervalle publié pour l'AUC regroupée ne peut pas être transféré à ce nouvel objet.

La référence de prévalence est constante \emph{au sein} de chaque année mais réestimée sur l'historique des années antérieures\cite{hgbcode}. Son AUC vaut donc 0,5 chaque année. Les effectifs observés et l'ordre de scores « 2025 au-dessus de 2024 » reconstituent exactement l'AUC regroupée 0,472568 ; c'est un contrôle algébrique de cohérence, pas une nouvelle vérification des scores individuels. Les précisions/rappels annuels au décile ont été calculés directement depuis les prédictions et sont rapportés ci-dessus ; ils ne se déduisent pas des seuls AUC, AP et Brier.

\subsection{CMS : effet des égalités sur la précision moyenne}
\label{sec:ties}
La réévaluation des prédictions individuelles fournit les douze valeurs annuelles ou regroupées de $\mathrm{AP}_{\mathrm{seuils}}$ (supplément S7). Pour la référence de prévalence, elles valent 0,140644 en 2024, 0,116118 en 2025 et 0,122069 sur les deux années regroupées. Les AP historiques correspondantes sont 0,148816, 0,118687 et 0,121754. Leur différence reflète le départage des scores égaux, pas une discrimination supplémentaire du score constant au sein d'une année.

Pour l'historique local, les AP par seuils sont 0,211180530 en 2024, 0,183722903 en 2025 et 0,196022698 dans le test regroupé. Elles coïncident avec les AP de rang. Il en va de même pour le modèle avec propriétaires. Les écarts propriétaires moins local restent donc $+0,001796527$, $-0,000175565$ et $+0,001249204$, respectivement. La convention ne change ni leur valeur ni leur signe sur ces prédictions; aucun intervalle historique n'est transféré à la nouvelle métrique.

Les frontières de décile des deux modèles appris sont uniques dans chaque année et dans le test regroupé : leur sélection ne dépend pas du départage. En revanche, la référence constante départage l'année entière, soit 8\,909 observations en 2024 et 10\,076 en 2025. Une sélection uniforme de même taille a pour précision attendue la prévalence annuelle, et non la précision historique issue de l'ordre date--CCN. Le supplément distingue cette référence uniforme d'un tirage limité aux ex æquo de frontière.

\subsection{Part D historique : extension inter-domaine et charge de sélection}

\begin{table}[!htbp]
\centering\small
\caption{Part D, test 2024, identité générique exacte. Les rappels sont à $K=10$. Rappel et MRR concernent les 1\,019 prescripteurs avec positifs ; P@10 utilise les 2\,000 prescripteurs admissibles. Le BPR macro/MRR est repris à six décimales du rapport. Sources : résumé et exécution compacte\cite{hgbresults,hgbpartd}.}
\label{tab:partd}
\input{tables/partd_fr.tex}
\end{table}

Dans le test historique v0.2, la popularité par spécialité obtient un rappel supérieur de 0,037071 au BPR relationnel ; l'IC complémentaire apparié par prescripteur pour BPR moins spécialité est $[-0,047179;-0,027336]$ (tableau~\ref{tab:ci}). Elle dépasse aussi la logistique de 0,049671 (tableau~\ref{tab:partd}). Ce classement est celui des méthodes alors comparées, avant la nouvelle sélection Q2.

L'ordre des heuristiques varie selon le seuil : à $K=5$, le recouvrement des historiques atteint 0,128561 contre 0,127465 pour la spécialité ; aucun intervalle n'est rapporté pour ce contraste. Cette comparaison conceptuelle bornée ne mesure pas l'effet propre de l'information relationnelle, faute de comparateur strictement local.

La cohorte de test compte 5\,476 liens positifs, répartis sur 1\,019 des 2\,000 prescripteurs. La popularité par spécialité retrouve 1\,193 liens parmi ses 20\,000 propositions. On obtient ainsi $1193/5476=0,217860$ de rappel, mais $1193/20000=0,059650$ de précision. Les 981 prescripteurs sans positif observé, soit 49,05\,\% de la cohorte, reçoivent eux aussi dix propositions\cite{hgbpartd}. Cette charge serait masquée par une analyse limitée au rappel des prescripteurs positifs. L'absence de lien publié ne permet toutefois pas de qualifier les autres propositions d'erreurs cliniques.


\subsection{Comparaisons Q2 réglées sur trois domaines}
\label{sec:q2-results}

Ces résultats sont issus d'ajustements Q2 nouveaux, distincts des tableaux historiques précédents. Les familles et configurations indiquées ont été choisies sur validation ; le test n'intervient pas dans leur sélection. Les tableaux résument les résultats des configurations sélectionnées sur les trois tâches\cite{hgbq2methods}. Les références à popularité globale, fréquence des voisins, popularité par spécialité et recouvrement restent fixes.

\begin{table}[!htbp]
\centering\small
\caption{Comparaison Q2 sur MAUDE : configurations des six familles réglables sélectionnées sur la validation 2023, puis rappel micro@10 moyen sur les tests 2024 et 2025. Les deux heuristiques fixes n'ont pas de sélection et leurs valeurs de validation ne sont pas rapportées ici. Lorsque trois répétitions sont requises, la valeur de test est leur moyenne avec les graines 103, 211 et 307 ; aucune graine n'est sélectionnée. Aucun nouvel intervalle n'est associé aux résultats.}
\label{tab:q2-maude}
% Derived from completed Q2 JSON; no new numerical experiment.
\begin{tabularx}{\linewidth}{@{}Xrrr@{}}
\toprule
Méthode & Val. 2023 & Test 2024 & Test 2025 \\
\midrule
Popularité globale & \textemdash & $0{,}182242$ & $0{,}186306$ \\
Fréquence des voisins & \textemdash & $0{,}189239$ & $0{,}194394$ \\
Logistique & $0{,}186762$ & $0{,}189738$ & $0{,}187003$ \\
HGB & $0{,}185853$ & $0{,}188739$ & $0{,}185051$ \\
Spectral & $0{,}176163$ & $0{,}188350$ & $0{,}181286$ \\
BPR & $0{,}227169$ & $0{,}230551$ & $0{,}239437$ \\
GraphSAGE mean & $0{,}190785$ & $0{,}204065$ & $0{,}211825$ \\
Contrôle none & $0{,}149167$ & $0{,}193403$ & $0{,}165388$ \\
\bottomrule
\end{tabularx}

\end{table}

En Q2, le BPR sélectionné dépasse les références fixes au test MAUDE : son rappel micro@10 est de 0,230551 en 2024 et 0,239437 en 2025, contre 0,189239 et 0,194394 pour les voisins (tableau~\ref{tab:q2-maude}). Le GraphSAGE \texttt{mean} atteint 0,204065 et 0,211825 ; \texttt{none} atteint 0,193403 et 0,165388. Ce classement ne réécrit pas le résultat historique v0.1 de MAUDE, où le GraphSAGE à trois époques restait sous la fréquence des voisins.

La configuration \texttt{none} retenue a passé les critères d'apprentissage sur sonde en validation puis au réajustement avant l'évaluation ; cela vaut pour ces ajustements Q2 et ne prouve ni une généralisation indépendante ni un bug historique de D1.

\begin{table}[!htbp]
\centering\small
\caption{Comparaison Q2 sur les inspections CMS des maisons de retraite : configurations sélectionnées selon l'AUC ROC de validation 2023, puis résultats des tests 2024, 2025 et regroupé. Les modèles à historique local et à ajout de variables de propriété sont appariés par famille ; les différences sont descriptives et aucun nouvel intervalle n'est calculé.}
\label{tab:q2-cms}
% Derived from completed Q2 JSON; no new numerical experiment.
\begin{tabularx}{\linewidth}{@{}Xrrrr@{}}
\toprule
Variables / modèle & Val. 2023 & 2024 & 2025 & 2024--25 \\
\midrule
Historique local / Logistique & $0{,}653494$ & $0{,}621082$ & $0{,}626874$ & $0{,}619416$ \\
Historique local / HGB & $0{,}580839$ & $0{,}608396$ & $0{,}630043$ & $0{,}617756$ \\
+ propriétaires / Logistique & $0{,}655366$ & $0{,}621832$ & $0{,}626598$ & $0{,}619574$ \\
+ propriétaires / HGB & $0{,}585064$ & $0{,}610319$ & $0{,}628614$ & $0{,}618872$ \\
\bottomrule
\end{tabularx}

\end{table}

Les écarts d'AUC regroupée entre ajout des variables de propriété et historique local restent faibles : $+0,000158$ pour la logistique et $+0,001115$ pour HGB (tableau~\ref{tab:q2-cms}). Leur signe varie entre 2024 et 2025 ; ces estimations ne suggèrent donc pas un bénéfice uniforme de l'ajout relationnel. Les intervalles historiques du modèle logistique v0.1 ne se transfèrent pas à ces nouveaux ajustements.

\begin{table}[!htbp]
\centering\small
\caption{Comparaison Q2 sur Medicare Part D : configurations sélectionnées sur validation 2023 et rappel micro@10 au test 2024. Le rappel a pour dénominateur tous les liens positifs observables ; les 2\,000 prescripteurs admissibles, y compris ceux sans positif, sont conservés dans la cohorte. Aucun échantillonnage de candidats ni nouvel intervalle n'est utilisé.}
\label{tab:q2-partd}
% Derived from completed Q2 JSON; no new numerical experiment.
\begin{tabularx}{\linewidth}{@{}Xrr@{}}
\toprule
Méthode & Val. 2023 & Test 2024 \\
\midrule
Popularité par spécialité & $0{,}220228$ & $0{,}217860$ \\
Recouvrement historique & $0{,}221954$ & $0{,}210373$ \\
Logistique & $0{,}194684$ & $0{,}184441$ \\
HGB & $0{,}195029$ & $0{,}197103$ \\
BPR & $0{,}235473$ & $0{,}236182$ \\
\bottomrule
\end{tabularx}

\end{table}

Pour Part D, le BPR Q2 atteint 0,236182 au test 2024, au-dessus de la popularité par spécialité (0,217860) et du recouvrement historique (0,2104) (tableau~\ref{tab:q2-partd}). Le résultat historique v0.2 avait le classement inverse entre spécialité (0,217860) et BPR (0,180789) : la nouvelle grille règle le BPR sur validation et change le classement, sans invalider le résultat historique ni isoler une cause architecturale. Les configurations, sensibilités par graine et ressources observées sont détaillées en S12 ; leurs coûts ne sont pas des budgets égaux.

\subsection{Prévision prospective de publication Part D}
\label{sec:q2-forecast}

Un exercice prospectif distinct scelle des scores en vue d'une publication officielle à venir de relations prescripteur--médicament de l'année de service 2025. Il s'agit de prédire la publication observée de liens, pas les prescriptions cliniques de 2025. Les six fichiers d'entrée 2019--2024 ont été recapturés ; une cohorte de 2\,000 prescripteurs excluant les populations de développement a reçu des scores des configurations Q2 choisies et des heuristiques fixes, sans accès aux labels cibles. Les scores sont scellés et ancrés publiquement alors que la source cible n'est pas répertoriée dans le catalogue observé\cite{hgbq2forecast}. La provenance et les limites de cette prévision figurent en S13.

\textbf{L'évaluation indépendante n'a pas encore été réalisée.} Le fichier officiel cible 2025 n'était pas listé dans les catalogues capturés, aucune métrique cible n'a été calculée et l'attestation humaine de non-consultation demeure inconnue. Le scellement ne vaut donc ni évaluation indépendante ni preuve d'indépendance scientifique.
\Needspace{10\baselineskip}
\subsection{Synthèse de l'incertitude}

Le tableau~\ref{tab:ci} conserve les contrastes historiques, sans agréger des tâches ni des métriques différentes. Les intervalles du rappel à dix propositions sont négatifs pour GraphSAGE à trois époques moins voisins sur MAUDE et pour BPR moins spécialité sur Part~D, et positifs pour voisins moins global sur MAUDE ; les trois intervalles CMS incluent zéro. Aucun de ces intervalles historiques ne décrit les comparaisons Q2. L'IC apparié GraphSAGE30--voisins appartient à l'analyse C distincte (tableau~\ref{tab:paired-maude}) ; il reste conditionnel à ses prédictions et ne couvre ni la sélection ni l'entraînement. Aucun nouvel IC n'est calculé pour les ajustements Q2.

\begin{table}[!htbp]
\centering\small
\caption{Contrastes et intervalles à 95\,\%. Les six premières lignes reprennent les intervalles publiés au v0.1 ; les deux dernières donnent les IC complémentaires, issus d'un bootstrap apparié post hoc sur prédictions gelées (1\,000 tirages, graine 20260929). Ces IC complémentaires sont conditionnels aux prédictions et ne couvrent pas la sélection ni l'entraînement. Sources : rapports v0.1/v0.2 et sorties indiquées au supplément S6\cite{hgbresults}.}
\label{tab:ci}
\begin{tabularx}{\linewidth}{@{}Xrr@{}}
\toprule
Contraste / mesure & Différence & Intervalle à 95\,\%\\
\midrule
MAUDE (v0.1) : GraphSAGE (3 époques) -- voisins, R@10 & $-0,019204$ & $[-0,023578;-0,014845]$\\
MAUDE (v0.1) : même contraste, macro R@10 & $-0,027883$ & $[-0,033630;-0,022316]$\\
MAUDE (v0.1) : même contraste, MRR & $-0,018949$ & $[-0,023026;-0,014850]$\\
CMS (v0.1) : propriétaires -- local, AUC & $+0,001318$ & $[-0,001458;+0,004329]$\\
CMS (v0.1) : même contraste, $\mathrm{AP}_{\mathrm{rang}}$ & $+0,001249$ & $[-0,001186;+0,003570]$\\
CMS (v0.1) : même contraste, Brier & $+0,000132$ & $[-0,000038;+0,000303]$\\
MAUDE (compl.) : voisins -- global, R@10 & $+0,007591$ & $[+0,005497;+0,009820]$\\
Part D (compl.) : BPR -- spécialité, R@10 & $-0,037071$ & $[-0,047179;-0,027336]$\\
\bottomrule
\end{tabularx}
\end{table}

\subsection{Contrôles synthétiques et tâche différée}

Un contrôle synthétique applique une règle de voisinage à un graphe de propriété : les étiquettes dépendent elles-mêmes de ce score. L'AUC du score relationnel est 0,503087 pour $\beta=0$ et 0,817456 pour $\beta=2,5$ (moyenne de trois graines)\cite{hgbcontrols}. Cette construction vérifie la présence d'un signal relationnel dans les données synthétiques ; elle n'évalue pas les modèles de santé ni un effet causal. Le mécanisme est détaillé en S1.4.

Un second contrôle exerce le noyau GraphSAGE sur un graphe synthétique à deux blocs (24 produits, 12 problèmes, 120 arêtes d'entraînement et 24 arêtes masquées ; tous les nœuds cibles sont connus)\cite{hgbcode}. Le rappel@1 des arêtes masquées est de 0,125000 sans apprentissage, 0,083333 après trois époques et 1,000000 après trente. Cet exemple montre que le noyau peut apprendre le signal construit ; il ne préjuge pas de la performance sur MAUDE ni de l'effet propre de l'agrégation. Les détails et contrôles de gradient sont en S8.

ClinicalTrials.gov/AACT reste séparé des trois tâches. La cible est la première publication des résultats dans les 365 jours suivant une origine située 0 à 90 jours après la fin primaire réelle ; il ne s'agit pas d'une étiquette de conformité juridique. Le contraste sponsors--variables locales donne $\Delta\mathrm{AP}=-0,048773$ en 2022 (IC $[-0,088028;-0,021689]$), et le contraste contexte hétérogène--sponsors $-0,108077$ en 2023 (IC $[-0,171780;-0,049169]$)\cite{hgb2026,hgbhistory}. Ces résultats concernent des origines distinctes et ne sont pas agrégés aux trois tâches.

L'historique des décisions indique que l'instabilité des gains a contribué au report de ClinicalTrials, tandis qu'aucune règle commune de succès relationnel antérieure aux résultats n'est établie pour l'ensemble des tâches. Cette asymétrie de sélection, documentée en S4, invite à interpréter les résultats comme des comparaisons rétrospectives plutôt que comme une validation confirmatoire.

\section{Discussion}

\subsection{Interprétation des résultats}

Le résultat historique à trois époques décrit une configuration, non une famille de modèles : l'analyse C sélectionne trente époques sur validation et obtient un rappel plus élevé au test. La sensibilité n'est pas monotone, puisque dix époques font moins bien que trois en validation. Le rapprochement de GraphSAGE30 avec les voisins conservés et son IC restent des résultats de C, distincts des nouveaux ajustements Q2. D1 demeure un diagnostic historique secondaire : sa perte presque plate sous le réglage imposé est conservée ; ni les critères d'apprentissage des nouveaux ajustements Q2 ni leurs résultats ne démontrent que D1 était un bug d'optimisation.

Les comparaisons Q2 changent certains classements sans les réécrire rétroactivement : sur MAUDE le BPR réglé dépasse les heuristiques fixes, tandis que le GraphSAGE \texttt{mean} et le contrôle \texttt{none} donnent des résultats différents selon l'année ; sur Part D, le BPR réglé dépasse la spécialité qui devançait le BPR historique. Ces différences ne démontrent pas un effet causal de l'architecture ou de l'agrégation : les ajustements, dimensions et capacités diffèrent, et \texttt{none} apprend toujours les liens via les vecteurs de nœuds et BPR. Sur CMS, les écarts liés aux propriétaires restent faibles et varient selon l'année et la famille. Les résultats sont propres aux méthodes et protocoles considérés ; ils ne fondent pas une hiérarchie générale des graphes.

Les mécanismes sous-jacents restent inconnus. Des agrégats simples peuvent capter une partie de l'information relationnelle, tandis qu'une architecture donnée peut ne pas l'exploiter efficacement ; ces explications demeurent des hypothèses. La comparaison CMS porte plus directement sur l'ajout de variables de propriété au même historique de base, alors que les méthodes MAUDE et Part D diffèrent aussi par leurs représentations et objectifs. Les résultats de RelBench rappellent que ces conclusions locales ne valent pas pour toutes les tâches relationnelles\cite{robinson2024relbench,gu2026relbench}.

\subsection{Trois limites de validité}

\paragraph{Validité interne.} Les périodes de test avaient déjà été consultées, y compris avant les nouvelles analyses MAUDE ; les choix faits sur validation ne rendent donc aucun de ces résultats indépendant. Les intervalles historiques disponibles sont conditionnels aux prédictions et ne couvrent pas la sélection ou l'entraînement ; aucun IC n'est calculé pour les comparaisons Q2. Les familles réglables ont le même nombre d'opportunités de sélection, mais pas des coûts de calcul égaux, d'autant que les exécutions ont partagé la machine. La prévision Part D reste scellée sans métrique cible ; l'évaluation indépendante n'a pas encore été réalisée et l'attestation humaine de non-consultation est inconnue.

\paragraph{Validité des mesures.} Les cibles proviennent d'observations publiées, sensibles au sous-signalement, au codage et aux changements de nomenclature. Une date d'événement ne prouve pas que tous les attributs étaient publiquement disponibles à cette date ; la disponibilité historique des fichiers Part D n'est pas vérifiée\cite{hgbresults}. Les limites de déclaration de la FDA empêchent aussi d'interpréter directement les scores comme un risque médical\cite{fda2026}.

\paragraph{Validité externe.} Ces trois tâches, les cohortes retenues et les périodes observées ne représentent pas l'ensemble des soins ou des bases relationnelles. L'admissibilité exige un historique et le classement se limite aux objets déjà présents dans le vocabulaire ; les scores ne mesurent ni les décisions prises par un utilisateur ni leurs conséquences.

\subsection{Bilan}

Pris ensemble, les résultats montrent que les classements dépendent de la tâche, des comparateurs et des configurations. Sur MAUDE, le résultat historique à trois époques ne caractérise pas toute la famille GraphSAGE : C retient trente époques, tandis que la comparaison Q2 place le BPR réglé devant les heuristiques fixes. D1 conserve son rôle diagnostique limité et sa perte presque plate. Sur CMS, l'incrément des variables de propriété reste faible et non uniforme. Sur Part D, les heuristiques dépassaient le BPR historique, alors que le BPR nouvellement réglé dépasse la spécialité au test Q2. Ces changements de classement sont exploratoires, sans nouvel intervalle ni interprétation causale.

Ni ces résultats ni la prévision scellée ne démontrent une supériorité générale des graphes ou des heuristiques, une utilité clinique ou une évaluation indépendante. Le benchmark fournit un cadre pour examiner ces comparaisons selon les tâches et leurs limites.


\section{Utilisation de la ressource et reproduction}
\label{sec:repro}
HealthGraphBench fournit une interface commune \texttt{load\_task}, des partitions \texttt{Split}, des prédictions \texttt{PredictionSet} et les méthodes \texttt{fit\_predict} et \texttt{evaluate}\cite{hgbinterface}. MAUDE et CMS nécessitent leurs sources externes. Pour Part D, le parcours actuel vérifie les six fichiers bruts, prépare les données puis ajuste réellement la méthode demandée et recalcule les métriques ; l'entrée préparée est également possible si elle est déclarée comme telle. L'appel ne consulte pas un classement \texttt{model-gate} conservé. L'univers des candidats reste défini par chaque tâche et les contrôles d'admissibilité correspondants.

La réutilisation technique en environnement wheel propre a parcouru données brutes, nouvelle préparation, ajustement, scores et évaluation publique. Elle a ajusté le modèle logistique SDK et le modèle \texttt{NeighborCosineTop50}, de similarité cosinus des voisins et implémenté hors du paquet. Ce modèle est appelé via l'interface publique sans ajout à la liste intégrée des modèles ; il utilise une incidence binaire prescripteur--médicament, retient au plus cinquante voisins et n'a ni apprentissage supervisé ni réglage. Ses scores ont été recalculés avec \texttt{PartDTask.evaluate}. Le rappel micro@10 obtenu par ce modèle est de 0,260960 en validation 2023 et de 0,258583 au test 2024 ; le témoin logistique atteint respectivement 0,178460 et 0,168188 (tableau~\ref{tab:q2-reuse}). Ces métriques sont exploratoires ; la logistique SDK conserve son algorithme historique à dix-huit époques et n'est pas le comparateur logistique LBFGS sélectionné dans la comparaison Q2.

\begin{table}[!htbp]
\centering\small
\caption{Témoin technique de réutilisation Part D : ajustements réalisés et métriques recalculées avec l'évaluateur public en validation 2023 et au test 2024. Le modèle cosinus est implémenté hors paquet et non réglé ; la logistique SDK n'est pas le comparateur Q2. Ces résultats ne constituent ni une comparaison réglée supplémentaire ni une validation humaine.}
\label{tab:q2-reuse}
% Derived from completed Q2 JSON; no new numerical experiment.
\begin{tabularx}{\linewidth}{@{}Xrr@{}}
\toprule
Méthode & Val. 2023 & Test 2024 \\
\midrule
Cosine top50 hors paquet & $0{,}260960$ & $0{,}258583$ \\
Logistique SDK (18 époques) & $0{,}178460$ & $0{,}168188$ \\
\bottomrule
\end{tabularx}

\end{table}

Cette démonstration technique, réalisée par l'assistant, établit un parcours logiciel, pas une réutilisation par un participant humain extérieur ni une évaluation indépendante. Aucun participant externe n'a réalisé le parcours ; l'évaluation indépendante liée à la prévision n'a pas encore eu lieu. Les preuves et limites sont décrites en S14.

Les sorties compactes historiques restent conditionnelles aux résultats conservés : les analyses C et D1 reposent sur des entraînements antérieurs, non relancés ici. S11 distingue leur calcul original du rejeu de résultats conservés. Ni ces vérifications ni le témoin technique Q2 n'établissent une disponibilité publique historique des champs ou une réplication indépendante.

\section{Conclusion}

HealthGraphBench documente trois tâches hétérogènes, sans estimer une valeur universelle des graphes. Les résultats historiques demeurent identifiables : MAUDE à trois époques favorise les voisins face à GraphSAGE, l'analyse C sélectionne trente époques et rapporte son propre contraste apparié, et D1 conserve une perte presque plate sous réglages imposés. Les nouvelles comparaisons Q2 changent le classement au sein des périodes exploratoires : le BPR réglé domine les heuristiques fixes sur MAUDE en rappel@10 ; sur Part D, le BPR Q2 dépasse la popularité par spécialité, à l'inverse du résultat historique ; sur CMS, l'apport des variables de propriété est faible et non uniforme. Le contrôle MAUDE sans propagation ne prouve pas un mécanisme causal ni un bug historique. Aucun nouvel intervalle n'est calculé pour Q2.

Les résultats ne montrent pas une supériorité générale des graphes ou des heuristiques ni une utilité clinique. La prévision prospective concerne la publication de relations Part D, non une issue clinique ; l'évaluation indépendante n'a pas encore été réalisée. L'ensemble reste un benchmark exploratoire dont les comparaisons doivent être lues avec leurs périodes, sélections et limites propres.
